% !PRINTED 487
\weberpage{487}
%<<CONT
\[
\psi_1(\alpha, y) = 0,
\]
so können wir wegen der Transitivität der Gruppe auf diese Gleichung eine Permutation anwenden, durch die $\alpha$ in eine beliebige Grösse der Reihe $A$ übergeht, wodurch $y$ ungeändert bleibt; es ist also jede Wurzel von $\psi(t, y)$ zugleich Wurzel von $\psi_1(t, y)$ und also sind beide Functionen identisch.

Es ist also die Gleichung $n^{\text{ten}}$ Grades $f(x) = 0$ imprimitiv in dem oben festgesetzten Sinne, d.\,h. eine Wurzel $\alpha$ dieser Gleichung kann betrachtet werden als Wurzel einer Gleichung $r^{\text{ten}}$ Grades, deren Coëfficienten von der Wurzel $y$ einer Gleichung $s^{\text{ten}}$ Grades abhängen.

Um zu beweisen, dass dann auch der Körper $\Omega(\alpha)$ (in dem Sinne des \S.\,144) imprimitiv ist, haben wir noch nachzuweisen, dass es in $\Omega(\alpha)$ Grössen giebt, die einer irreducibeln Gleichung $s^{\text{ten}}$ Grades genügen.

Wir werden nachweisen, dass $y$ selbst eine solche Grösse ist, dass also $y$ rational durch $\alpha$ allein darstellbar ist.

Dies ist aber sehr einfach zu schliessen. Nach (11) ist
\weq{12}{\psi(\alpha, y) = 0,}
während $\psi(\alpha, y_1), \ldots \psi(\alpha, y_{s-1})$ von Null verschieden sind. Denn wäre etwa $\psi(\alpha, y_1) = 0$, so wäre $\alpha$ eine Wurzel der Gleichung $r^{\text{ten}}$ Grades $\psi(t, y_1) = 0$, deren Wurzeln ja die von $\alpha$ verschiedenen Grössen $\beta, \beta_1 \ldots \beta_{r-1}$ sind. Es haben also die beiden Gleichungen
\weq{13}{\psi(\alpha, u) = 0, \quad \varphi(u) = 0}
eine und nur eine gemeinsame Wurzel $u = y$. Der grösste gemeinschaftliche Theiler dieser beiden Functionen ist daher in Bezug auf $u$ linear, und er giebt $y$ rational durch die Coëfficienten von $\varphi(u)$ und $\psi(\alpha, u)$, \gesperrt{also rational durch} $\alpha$. Damit ist die Umkehrung von 1. bewiesen:

3. \gesperrt{Ein primitiver Körper hat eine primitive Gruppe.}

\begin{center}
\rule{3cm}{0.4pt}
\end{center}
